\documentclass[11pt,a4paper]{article}
\usepackage[utf8]{inputenc}
\usepackage[T1]{fontenc}
\usepackage[ngerman]{babel}
\usepackage[margin=0.75in]{geometry}
\usepackage{enumitem}
\usepackage{titlesec}
\usepackage{xcolor}
\usepackage{hyperref}
\usepackage{lmodern}
\input{glyphtounicode}
\pdfgentounicode=1
\definecolor{primarycolor}{RGB}{33,37,41}
\definecolor{secondarycolor}{RGB}{0,102,204}
\definecolor{accentcolor}{RGB}{108,117,125}
\hypersetup{colorlinks=true,linkcolor=secondarycolor,urlcolor=secondarycolor,citecolor=secondarycolor}
\titleformat{\section}{\large\bfseries\color{primarycolor}}{}{0em}{}
\titlespacing*{\section}{0pt}{0.9em}{0.4em}
\newcommand{\cvsection}[1]{\section*{#1}}
\newcommand{\cvrole}[2]{\noindent\textbf{#1}\\#2\par}
\pagestyle{empty}
\setlength{\parindent}{0pt}
\setlength{\parskip}{0.25em}
\setlist[itemize]{leftmargin=*,itemsep=0.15em,topsep=0.25em,parsep=0pt}
\begin{document}
{\LARGE\bfseries\color{primarycolor} KATARZYNA WIŚNIEWSKA}\par
{\large\color{accentcolor} Pflegehelferin}\par
Habsburgerstraße 41, 79104 Freiburg im Breisgau\\
k.wisniewska@example.com | +49 761 555 0412

\cvsection{PROFIL}
Pflegehelferin mit Erfahrung in der stationären Altenpflege seit April 2020 und vorheriger Tätigkeit als Krankenschwester in Polen. Praxis in der Grundpflege, Qualifikation Behandlungspflege LG1 und sicherer Umgang mit Vivendi PD. Ruhig, zuverlässig und geduldig im Umgang mit älteren Menschen und Menschen mit Demenz. Deutschkenntnisse auf Niveau B2.

\cvsection{BERUFSERFAHRUNG}
\cvrole{Pflegehelferin | Seniorenzentrum Am Lindenhof, Freiburg}{04/2020 – heute}
\begin{itemize}
\item Grundpflege für 10 bis 12 Bewohnerinnen und Bewohner pro Schicht: Körperpflege, Mobilisation, Lagerung und Hilfe beim Essen.
\item Behandlungspflege nach Leistungsgruppe 1 (LG1): Blutdruck- und Blutzuckermessung, Kompressionsstrümpfe und Medikamente nach Plan reichen.
\item Pflegedokumentation in Vivendi PD.
\item Erfahrung in Früh-, Spät- und Wochenenddiensten sowie im Umgang mit Menschen mit Demenz.
\item Einarbeitung von bisher sieben neuen Pflegehelferinnen und Praktikanten während der ersten zwei Wochen nach dem Einarbeitungsplan der Einrichtung.
\end{itemize}

\cvrole{Aushilfe Verkauf | Bäckerei Schwarzwaldkorn, Freiburg}{09/2019 – 03/2020}
\begin{itemize}
\item Verkauf, Kasse und Reinigung.
\end{itemize}

\cvrole{Krankenschwester (Pielęgniarka) | Wojewódzki Szpital Specjalistyczny, Opole, Polen}{10/2014 – 06/2019}
\begin{itemize}
\item Grund- und Behandlungspflege auf einer Station für Innere Medizin mit 32 Betten.
\item Infusionen, Wundversorgung und Vorbereitung von Visiten.
\item Verbandswechsel bei Dekubitus und diabetischem Fuß unter Aufsicht der Stationsärztin; diese Erfahrung wurde in Polen erworben.
\end{itemize}

\newpage
\cvsection{AUSGEWÄHLTE PRAXISPROJEKTE}
\textbf{Sturzprophylaxe im Wohnbereich 2 | 2023}\\
Seniorenzentrum Am Lindenhof
\begin{itemize}
\item Gemeinsam mit der Wohnbereichsleitung Sturzprotokolle aus sechs Monaten ausgewertet; Nachtlichter und feste Toilettengänge für neun Bewohner eingeführt.
\item Im folgenden Halbjahr sank die Zahl der Stürze im Wohnbereich von 14 auf 8.
\end{itemize}

\textbf{Nachmittagsgruppe für Menschen mit Demenz | 2022 – heute}\\
Seniorenzentrum Am Lindenhof
\begin{itemize}
\item Wöchentliche Singgruppe mit polnischen und deutschen Volksliedern für sechs bis acht Bewohner geleitet.
\item Positive Rückmeldungen von Angehörigen im Heimbeirat.
\end{itemize}

\cvsection{AUSBILDUNG UND ANERKENNUNG}
\textbf{Bachelor Pielęgniarstwo (Krankenpflege)}\\
Medizinische Universität Opole | 2014

\textbf{Anerkennungsverfahren als Pflegefachfrau}\\
Regierungspräsidium Stuttgart | laufend
\begin{itemize}
\item Defizitbescheid liegt vor; Kenntnisprüfung für März 2027 vorgesehen.
\item Derzeit in Deutschland als Pflegehelferin tätig; Anerkennung als Pflegefachfrau noch nicht abgeschlossen.
\end{itemize}

\cvsection{WEITERBILDUNG}
\begin{itemize}
\item Erste-Hilfe-Kurs | 2024
\item Expertenstandard Sturzprophylaxe, interne Schulung | 2023
\item Behandlungspflege LG1, Basiskurs | 2021
\end{itemize}

\cvsection{KENNTNISSE UND SPRACHEN}
\begin{itemize}
\item Pflegedokumentation: Vivendi PD, sicherer Umgang.
\item Polnisch: Muttersprache.
\item Deutsch: B2, telc Deutsch B1-B2 Pflege, Zertifikat 2021.
\item Englisch: Grundkenntnisse.
\end{itemize}

\cvsection{FÜHRERSCHEIN}
Klasse B noch nicht vorhanden; Theorieprüfung bestanden, praktische Prüfung für Januar 2027 vorgesehen.
\end{document}